\scanpage{16}\mspage{15}
par contre la relation $\ll$ n'a pas d'interprétation simple directe, ni en termes des \inserted{\unsure{déployements|déploiements}} $\widetilde F\subset M$, ni en termes des \unsure{multidéployement|multidéployements|multidéploiement|multidéploiements} $\mathrm{Fig}(F)$; ainsi $G\ll F$ n'implique nullement $\mathrm{Fig}_{M}(G)\ll\mathrm{Fig}_{M}(F)$. Aussi l'application ($**$) et la notation $=\mathrm{Fig}_{M}(F)$ sont-ils finalement un peu bidon !

En termes de $M$, la structure est donnée par
\[
\boxed{\begin{aligned}
&\unsure{\mathcal{F}|\mathscr{F}}\subset\unsure{\mathcal{P}|\mathrm{P}}(M)\quad\text{ens de parties de }M\\
&\ll\quad\text{relation d'ordre dans }\mathcal{F},\ \text{\uline{impliquée}}\\
&\qquad\text{par l'inclusion }(\deleted{\unsure{F}}\text{notée }F\leq G)
\end{aligned}}
\tag*{(1.33)}
\]
Ça a l'air moins sympathique du point de vue simplicité des données et d'autant plus qu'il n'y a pas \inserted{à première vue} d'interprétation directe de $\ll$ en termes de \deleted{\unsure{notions}\ \unsure{de structure|de théorie}} \inserted{notions ensemblistes} familières. Mais c'est ici que vont servir les (grandes) \unsure{ombres} !

En termes de la \inserted{première} donnée (1.33\unsure{.}\unsure{)|]} \margin{(et avant d'introduire $\ll$)}, le \deleted{\unsure{tout}} premier axiome est \deleted{\illegible{1}} le fait que $\mathcal{F}$ est \deleted{une figure} \inserted{une famille d'}ens. \inserted{\uuline{pré}admissibles} \deleted{\unsure{dans}} $M$ de
% slanted
\margin{attention au ``\uuline{pré}'' !}
\deleted{\unsure{Axiome}} support $M$, \deleted{\unsure{d'où}} donc la donnée de
% slanted
\margin{\deleted{\unsure{un ens.|un ensemble} de parties \unsure{fermées} de $\unsure{F|\mathcal{F}|\lbrace}$}}
$\mathcal{F}$ \deleted{\unsure{revient à|revient} la donnée d'} \inserted{implique} une relation de préordre dans $M$ (ce sera $\leq$)\unsure{$\big/\!6\!\big/$} -- lequel sera
